\section{Fejlett kliensoldali web (AJAX, SPA-k, keretrendszerek)}
A 2000-es évek közepétől a webalkalmazások egyre inkább „alkalmazásszerűvé” váltak \cite{mdnClientFrameworks}. Az AJAX (a kifejezést 2005-ben vezették be) lehetővé tette, hogy a böngészőben futó JavaScript aszinkron módon kérjen le adatokat \cite{garrett2005ajax,mdnAjax}. Az olyan keretrendszerek, mint a jQuery (2006), Angular (2010), React (2013), Vue stb., lehetővé tették az egyoldalas alkalmazások (SPA-k) létrehozását, ahol a tartalom dinamikusan töltődik be kliensoldalon \cite{mdnSpa,mdnClientFrameworks}. Ezek az alkalmazások gyakran REST vagy GraphQL API-kkal kommunikálnak \cite{fielding2000rest,graphql2021spec}.

A HTML5 (W3C ajánlás, 2014 október) szabványosított új API-kat (Canvas, WebSockets, Web Storage, videó/hang támogatás) a fejlettebb funkciók érdekében \cite{w3c2014html5,fette2011websocket,w3c2013webstorage}. Ennek eredményeként a modern weboldalak képesek összetett logikát futtatni a böngészőben, offline gyorsítótárazást és push értesítéseket használni \cite{mdnPwaOfflineBackground}. A „statikus weboldal” és a „webalkalmazás” közötti különbség elmosódott, mivel még a tartalomközpontú oldalak is átvették a kliensoldali renderelést \cite{webdevClientSideRendering,mdnClientFrameworks}.
